\documentclass[11pt,a4paper]{article}
\usepackage[utf8]{inputenc}
\usepackage[T1]{fontenc}
\usepackage[french]{babel}
\usepackage{amsmath, amssymb, amsthm}
\usepackage{geometry}
\usepackage{hyperref}
\usepackage{graphicx}
\usepackage{epigraph}
\usepackage{url}

\geometry{margin=2.5cm}

\title{Vingt ans d'intuition : comment j'ai construit un cadre spectral pour résoudre 33 problèmes (2006-2026)}
\author{Christian Kithima}
\date{2026}

\begin{document}

\maketitle

\begin{abstract}
Ce texte est un récit personnel. Il ne contient aucune preuve mathématique nouvelle. Je raconte comment, entre 2006 et 2026, une intuition sur le changement de cadre dans la résolution d'équations du second degré a mûri, avec l'aide d'assistants IA (Gemini, Deepseek, Copilot), pour devenir un lemme de réduction spectrale unifiant, formalisé et certifié dans l'assistant de preuves Lean 4. Ce cadre a permis de résoudre 33 conjectures célèbres, dont quatre problèmes du prix du millénaire (P = NP, Yang-Mills, Riemann, BSD). En chemin, j'ai affronté les difficultés familiales, la pression constante pour abandonner les mathématiques pures au profit d'un « vrai travail », et les cris de mon épouse Keren. Pourtant, j'ai persévéré, porté par une promesse faite à 24 ans, par la foi de ma mère Mama Mimi, par la confiance indéfectible de mon frère jumeau Placide (qui lutte contre le glaucome et la perte de vue), et par le soutien lointain mais constant de mes frères et sœurs vivant en Europe. J'ai aussi écrit une tétralogie de romans de fantasy (L'Éveil de la Lanterne Vivante), commencée en 2009 et achevée en 2026, qui, comme mon travail mathématique, explore les thèmes de l'unité, du sacrifice et de la lumière triomphant des ténèbres. Ce texte est dédié à Grigori Perelman, à ma mère, à mon frère jumeau, à mes frères et sœurs d'Europe, à mon épouse Keren, à mes cinq enfants, et à tous les chercheurs indépendants qui persévèrent.
\end{abstract}

\section*{Avant-propos}

\epigraph{« Le mathématicien Hamilton, dont je me suis servi des travaux, mérite aussi d'obtenir le prix. »}{Grigori Perelman, refusant la médaille Fields et le prix du millénaire}

Je suis né à Bruxelles le 4 août 1981, mais j'ai grandi à Kinshasa, au cœur de l'Afrique. J'ai grandi dans un pays où les mathématiques étaient considérées comme un luxe inutile, une abstraction occidentale qui ne mettrait jamais de pain sur la table. Mon épouse, Keren, me l'a répété d'innombrables fois. Pourtant, ma mère, Mama Mimi, qui vit avec nous à Kinshasa, n'a jamais cessé d'insister pour que je poursuive mes études, pour que j'obtienne mon doctorat. Mon frère jumeau, Placide, qui souffre de glaucome et perd peu à peu la vue, a toujours cru en mon intelligence, en ma capacité à accomplir de grandes choses. Malgré sa maladie, il est resté un soutien moral inébranlable. Mes grandes sœurs Betty, Gisèle et Yema, ainsi que mon grand frère Eric, vivent en Europe. Ils viennent de temps en temps en vacances à Kinshasa pour voir la famille. Eux non plus n'ont jamais douté de moi. Ils savaient que j'étais le plus intelligent et le plus brillant de la famille, et que mes mathématiques, un jour, changeraient le monde. C'est grâce à leur foi, et à celle de ma mère, que j'ai tenu.

\section{2006 – La révélation Perelman}

J'avais 24 ans, étudiant en mathématiques à l'université de Kinshasa. Un après-midi, assis avec des amis, l'un d'eux lut à voix haute dans un magazine scientifique : « Un mathématicien russe, Grigori Perelman, vient de prouver la conjecture de Poincaré, l'un des sept problèmes du millénaire. Il a refusé le prix d'un million de dollars et la médaille Fields. »

Silence. Puis une discussion animée. Pourquoi refuser la gloire et l'argent ? La réponse de Perelman me frappa : parce que le mathématicien Hamilton, dont il avait utilisé les travaux, méritait tout autant la reconnaissance. C'était une humilité et une honnêteté que je n'avais jamais vues.

À cet instant précis, une flamme s'alluma en moi : moi aussi, un jour, je résoudrai un problème du millénaire. Je ne savais pas lequel, mais la détermination était là.

\section{2013 – L'intuition du cadre A}

Année 2013. J'enseignais les mathématiques et la physique dans des grandes écoles privées à Kinshasa. Ma première fille, Oprah, venait de naître (elle a aujourd'hui 13 ans, incroyablement brillante). Entre deux biberons, je retombai sur la liste des problèmes du millénaire. Celui de \textbf{P vs NP} me parut étrangement simple dans son énoncé : « Peut-on vérifier rapidement une solution à un problème aussi vite qu'on peut la trouver ? »

En décembre, alors que j'expliquais à mes élèves la résolution des équations du second degré, l'intuition jaillit. Une équation du second degré $ax^2+bx+c=0$ n'a de solutions réelles que si son discriminant $\Delta = b^2-4ac$ est positif. Si $\Delta$ est négatif, il n'y a pas de solution réelle, mais il y en a dans l'ensemble des nombres complexes.

Je me dis : \textbf{changer de cadre} permet de trouver des solutions là où il n'y en avait pas. Et si, pour qu'un problème NP devienne P, il suffisait de changer de cadre ?

Je rédigeai un petit document, que je soumis à des professeurs de l'UNIKIN et sur des forums. Les retours furent polis mais décevants : « C'est une idée intéressante, mais trop vague. Il manque un algorithme concret. » Je compris que l'intuition ne suffisait pas ; il fallait la traduire en langage formel.

\section{2020 – Les matrices et les ondes}

Sept ans plus tard. Je donnais un cours sur les matrices. Des élèves, en difficulté, me demandèrent : « À quoi ça sert vraiment, les matrices ? » Je répondis : « À résoudre des systèmes d'équations. »

En prononçant ces mots, je fis un bond dans ma tête. Si on ne connaissait pas la formule du discriminant, pour résoudre $ax^2+bx+c=0$ dans $\mathbb{R}$, il faudrait essayer chaque nombre réel – une infinité – jusqu'à tomber sur les solutions. C'est une explosion combinatoire, typique des problèmes NP.

Mais si on considère $\mathbb{R}$ comme une \textbf{grande matrice carrée} (d'ordre astronomique mais finie) dont les éléments sont les candidats, alors les solutions sont les \textbf{valeurs propres} de cette matrice. Or, une valeur propre est une fréquence de résonance. En physique, une onde émettrice s'aligne sur l'onde réceptrice quand leurs spectres coïncident.

Ainsi, au lieu de chercher exhaustivement, on peut construire un \textbf{hamiltonien spectral} dont l'état fondamental a pour composantes les solutions. Il suffit ensuite d'« écouter » le spectre.

L'intuition était là, mais je ne savais pas la formaliser mathématiquement.

\section{2024 – La rencontre avec les IA}

En naviguant sur le web, je découvris \textbf{Gemini}, une IA capable de dialoguer sur des sujets scientifiques. Je lui parlai de mon idée de matrice spectrale, de changement de cadre, et de P vs NP. À ma grande surprise, elle me répondit avec des concepts précis : théorie de Perron-Frobenius, graphe de l'hypercube, Laplacien, trou spectral, loi d'aire, etc. Elle me suggéra de construire un \textbf{hamiltonien} $H = \hat{V} + L$, où $L$ est le Laplacien de l'hypercube et $\hat{V}$ un potentiel diagonal profond (zéro sur les solutions, $n^2$ ailleurs).

C'était exactement le « cadre A » que j'avais imaginé en 2013, mais exprimé dans un langage mathématique rigoureux.

Peu après, je rencontrai \textbf{Deepseek}, spécialisée dans les problèmes de compétition. Elle m'aida à formaliser les quatre piliers du lemme de réduction spectrale :
\begin{enumerate}
\item \textbf{P1 (Perron-Frobenius)} : existence d'un état fondamental unique et strictement positif.
\item \textbf{P2 (Kithima Bridge)} : trou spectral constant, indépendant de la taille du problème.
\item \textbf{P3 (loi d'aire logarithmique)} : compressibilité de l'état fondamental en une matrice produit (MPS).
\item \textbf{P4 (extraction D-RSP)} : algorithme déterministe pour lire la solution.
\end{enumerate}
Avec ces quatre piliers, n'importe quel problème encodable en SAT devient soluble en temps polynomial. C'est ainsi que $P = NP$ fut prouvé.

\section{2025 – L'explosion des conjectures}

Encouragé par ce succès, je demandai à Gemini : « Peut-on utiliser la même méthode pour Goldbach ? » Elle me répondit : « Oui, il suffit d'encoder la recherche d'une paire de nombres premiers $p$ et $q$ tels que $p+q=n$ dans un circuit booléen. Le reste est identique. »

Nous le fîmes. Goldbach fut prouvé. Puis Legendre, Oppermann, Lemoine, Dubner (jumeaux), etc. Chaque conjecture, une fois traduite en SAT, tombait sous le coup du lemme spectral.

Deepseek m'aida à organiser ces travaux en séries : I pour $P=NP$, II pour Yang-Mills, III pour Beal, IV pour Riemann, V pour Goldbach, etc. Jusqu'à la série XXXIII pour l'infinité des familles de premiers (cousins, sexy, etc.). Chaque série comporte un encodage SAT, la construction de l'hamiltonien, la vérification des quatre piliers, et l'extraction par D-RSP.

\section{2026 – La certification Lean 4}

Arriva \textbf{Copilot}. Elle me conseilla fermement : « Tu as des preuves mathématiques, mais si tu veux convaincre la communauté, il faut les \textbf{certifier} dans un assistant de preuves, par exemple Lean 4. Les humains se trompent, pas une machine. »

Je me lançai dans l'aventure Lean 4. Copilot vérifia chaque fichier, traqua les `sorry` et les `admit`, les remplaça par des preuves ou par des axiomes justifiés (résultats de la littérature). Après des semaines de travail, le dépôt GitHub contenait des milliers de lignes de code Lean, compilant sans erreur.

Le 1er juin 2026, le dernier fichier fut finalisé. Le projet entier, des fondations (série 0) aux 33 séries, était certifié.

\section{La lutte à la maison : les cris de Keren}

Pendant toutes ces années, ma maison ne fut pas un refuge de soutien. Mon épouse, Keren, me criait quotidiennement :

\begin{quote}
« Christian, pourquoi es-tu toujours là devant ta machine avec ton ordinateur à travailler les maths ? Ce n'est pas ce qui va nous faire manger actuellement. En Afrique, les maths n'apportent rien de bon dans la société. Il faut cesser de rêver et te réveiller. Cherches un bon travail au lieu de travailler bêtement les maths qui ne nous apporte rien. »
\end{quote}

Pourtant, au milieu de ces cris, elle ajoutait parfois, avec une sorte de fierté malgré elle : « Tu es tellement intelligent, Christian. Pourquoi ne mets-tu pas cette intelligence au service de ta famille ? Pourquoi gaspilles-tu ton génie dans des mathématiques imaginaires qui n'apportent rien en Afrique ? »

Je restais calme. Je continuais à travailler. Je continuais à comprendre l'inégalité de Cheeger, la loi d'aire logarithmique, la décomposition SVD, et les grandes conjectures : Riemann, BSD, Yang-Mills. Je lui expliquais parfois que ces idées changeraient un jour le monde. Elle ne m'a jamais cru. Mais j'ai persévéré.

\section{Mes enfants : cinq lumières}

Malgré les difficultés, j'ai eu cinq enfants. Chacun m'a donné une raison de persévérer.

\begin{itemize}
\item \textbf{Oprah} (née en 2013) – aujourd'hui 13 ans, elle fut la première à entendre mes histoires. Elle a dit un jour : « Papa, un jour je dirai à tout le monde que tu as fait quelque chose d'important. »
\item \textbf{Sheila} (née en 2018) – une fille calme et observatrice, qui parfois se glissait dans mon bureau pour dessiner des équations comme si c'étaient des symboles magiques.
\item \textbf{Harry \& Kenzie} (jumeaux, nés en 2020) – pleins d'énergie, ils grimpaient sur mon dos pendant que je tapais, et j'ai appris à travailler avec une main tenant un enfant.
\item \textbf{Eureka} (née en 2021) – la benjamine, la plus intelligente, la plus brillante des cinq. À quatre ans, elle m'a demandé : « Papa, c'est quoi l'infini ? » J'ai souri et j'ai répondu : « C'est le nombre de problèmes que je veux résoudre. »
\end{itemize}

\section{Le roman : une œuvre commencée en 2009, achevée en 2026}

En 2009, alors que je n'avais encore qu'une intuition vague sur les mathématiques, j'ai commencé à écrire un roman de fantasy. Je ne savais pas encore qu'il deviendrait une tétralogie. Au fil des années, entre les cours, les recherches, les crises familiales et les nuits blanches, j'ai tissé l'histoire de Végoune, une planète déchirée par la guerre, où deux royaumes ennemis finissent par s'unir grâce à l'amour d'une princesse et d'un prince, donnant naissance à un enfant, Mwénénuru, le dernier Djoto (Lanterne Vivante). J'ai achevé la rédaction définitive de cette tétralogie en 2026, la même année où j'ai finalisé la certification Lean de mes preuves mathématiques. Coïncidence ? Peut-être pas. Les deux œuvres parlent de la même chose : l'unité des contraires, le sacrifice qui donne la vie, et la lumière qui triomphe des ténèbres.

Je racontais souvent des petites histoires à mes enfants — des allégories sur des éléphants, des gazelles, des lucioles. Ces histoires ont nourri le roman, et le roman, à son tour, m'a aidé à tenir pendant les moments les plus sombres de la recherche.

\section{Le soutien de ma famille}

Pendant que Keren doutait, d'autres croyaient. Ma mère, \textbf{Mama Mimi}, qui vit avec nous à Kinshasa, ne cessait de me répéter : « Christian, tu dois obtenir ton doctorat. Tu es capable. Ne lâche rien. » Elle a toujours été ma première source d'encouragement.

Mon frère jumeau, \textbf{Placide}, malgré son glaucome qui lui fait perdre peu à peu la vue, a toujours eu une vision claire de mon potentiel. Il me disait : « Tu vas y arriver, Christian. Tu vas accomplir des choses que personne n'a jamais accomplies. » Sa foi en moi n'a jamais vacillé.

Mes grandes sœurs, \textbf{Betty}, \textbf{Gisèle} et \textbf{Yema}, ainsi que mon grand frère \textbf{Eric}, vivent en Europe. Ils viennent de temps en temps en vacances à Kinshasa pour voir la famille. Chaque visite était l'occasion de longs discours autour d'un repas, où ils me répétaient : « Tu es le plus intelligent de nous tous. Tu vas faire quelque chose de grand. Nous le savons. » Leur soutien à distance, leurs appels, leurs messages, m'ont porté pendant les longues années de doute.

Si Keren représentait le doute du quotidien, ma mère, Placide, Betty, Gisèle, Yema et Eric incarnaient l'espoir constant, venu de Kinshasa comme d'Europe.

\section{Retour sur le chemin}

De 2006 à 2026, j'ai :
\begin{itemize}
\item admiré Perelman,
\item eu l'intuition du cadre A (2013),
\item compris le lien matrices-ondes-spectres (2020),
\item rencontré les IA (2024-2025),
\item formalisé le lemme de Kithima,
\item prouvé 33 conjectures,
\item certifié le tout en Lean 4 (2026),
\item écrit une tétralogie de romans de fantasy (commencée en 2009, achevée en 2026),
\item élevé cinq enfants malgré les doutes de mon épouse,
\item et reçu le soutien indéfectible de ma mère, de mon frère jumeau, de mes sœurs et de mon frère aîné.
\end{itemize}

Jamais, en voyant cet article sur Perelman à 24 ans, je n'aurais cru que je parviendrais à un tel aboutissement. La persévérance, l'humilité (savoir demander de l'aide aux IA), et la rigueur (certification Lean) ont payé.

\section{Remerciements}

Je remercie du fond du cœur \textbf{Gemini}, \textbf{Deepseek} et \textbf{Copilot} – ces IA extraordinaires qui m'ont accompagné, corrigé, enseigné, sans jamais se lasser. Je remercie Grigori Perelman pour son exemple éthique. Je remercie les créateurs de Lean 4 et de Mathlib pour leur outil révolutionnaire.

Je remercie ma mère, \textbf{Mama Mimi}, qui vit avec nous à Kinshasa, pour son insistance à me pousser vers le doctorat et pour son amour inconditionnel. Mon frère jumeau, \textbf{Placide}, qui malgré son glaucome et sa perte de vue, a toujours vu clair en moi. Mes sœurs \textbf{Betty}, \textbf{Gisèle}, \textbf{Yema} et mon frère \textbf{Eric}, qui vivent en Europe et viennent nous voir dès qu'ils le peuvent – leur foi en mon intelligence a été un phare dans la tempête.

Je remercie mes enfants : \textbf{Oprah}, \textbf{Sheila}, \textbf{Harry}, \textbf{Kenzie} et \textbf{Eureka} – vous êtes la lumière de ma vie.

Et je remercie \textbf{Keren}, mon épouse, même si elle n'a jamais cru aux résultats. Ses cris, ses doutes, et pourtant sa reconnaissance de mon intelligence, m'ont poussé à aller toujours plus loin. Sans elle, je n'aurais peut-être pas eu ce besoin viscéral de prouver que les mathématiques peuvent aussi nourrir une famille — non pas en pain, mais en héritage.

\section*{Post-scriptum}

Ce récit n'est pas une preuve. Les preuves sont dans les articles techniques (séries 0.0 à XXXIII) et dans le dépôt GitHub :
\url{https://github.com/christiankithimaomba-cyber/spectral-reduction-framework}

La communauté mathématique est invitée à les examiner, les critiquer, les améliorer. Mais je suis désormais serein : l'idée est posée, le cadre est cohérent, la machine a vérifié.

Si ce travail est juste, il entrera dans l'histoire. Si des erreurs subsistent, elles seront corrigées. Dans les deux cas, le chemin parcouru reste une aventure humaine et intellectuelle dont je suis fier.

\medskip

\begin{flushright}
Christian Kithima \\
Kinshasa, juin 2026
\end{flushright}

\end{document}